%%=============================================================================
%% Requirementsanalyse en Systeemarchitectuur
%%=============================================================================

\chapter{\IfLanguageName{dutch}{Requirementsanalyse en Systeemarchitectuur}{Requirements Analysis and System Architecture}}%
\label{ch:requirementsanalyse}

%% TODO: In dit hoofstuk geef je een korte toelichting over hoe je te werk bent
%% gegaan. Verdeel je onderzoek in grote fasen, en licht in elke fase toe wat
%% de doelstelling was, welke deliverables daar uit gekomen zijn, en welke
%% onderzoeksmethoden je daarbij toegepast hebt. Verantwoord waarom je
%% op deze manier te werk gegaan bent.
%% 
%% Voorbeelden van zulke fasen zijn: literatuurstudie, opstellen van een
%% requirements-analyse, opstellen long-list (bij vergelijkende studie),
%% selectie van geschikte tools (bij vergelijkende studie, "short-list"),
%% opzetten testopstelling/PoC, uitvoeren testen en verzamelen
%% van resultaten, analyse van resultaten, ...
%%
%% !!!!! LET OP !!!!!
%%
%% Het is uitdrukkelijk NIET de bedoeling dat je het grootste deel van de corpus
%% van je bachelorproef in dit hoofstuk verwerkt! Dit hoofdstuk is eerder een
%% kort overzicht van je plan van aanpak.
%%
%% Maak voor elke fase (behalve het literatuuronderzoek) een NIEUW HOOFDSTUK aan
%% en geef het een gepaste titel.

Dit hoofdstuk vertaalt de theoretische kaders en wetmatigheden uit de literatuurstudie 
naar concrete, meetbare functionele en niet-functionele systeemeisen. Het doel van deze 
requirementsanalyse is het vastleggen van een objectief toetsingskader waaraan de 
fysieke proof-of-concept (PoC) in een latere fase wordt gevalideerd. Vervolgens wordt 
op basis van deze eisen de logische en fysieke systeemarchitectuur ontworpen, waarbij 
de geselecteerde netwerkprotocollen, transportmechanismen en spectrumkeuzes 
structureel worden gedefinieerd.

\section{Netwerktechnische Requirements (RFC-Stijl)}
\label{sec:req-netwerk}

Om te waarborgen dat stateful applicatielaagprotocollen stabiel kunnen functioneren over 
een stochastisch radiomedium, moeten de prestatie-indicatoren van de netwerklaag en 
transportlaag binnen strikte limieten blijven. Op basis van de netwerkarchitecturen 
beschreven door \textcite{Tanenbaum2021} zijn de volgende niet-functionele 
netwerkeisen gedefinieerd.

\subsection{Minimale bandbreedte (REQ-01)}
Gezien het theorema van \textcite{Shannon1948}, waarin de kanaalcapaciteit inherent 
gelimiteerd wordt door de bandbreedte ($W$) en de signaal-ruisverhouding ($S/N$), moet 
de opstelling een minimale netto bitsnelheid (*goodput*) garanderen op de transportlaag. 
\begin{itemize}
    \item \textbf{Eis:} De effectieve doorvoersnelheid voor bestandsoverdracht moet minimaal 
    $1\text{ Mbit/s}$ bedragen onder suboptimale condities (smalle band/lange afstand) en 
    minimaal $10\text{ Mbit/s}$ onder optimale *Line-of-Sight*-condities.
    \item \textbf{Verantwoording:} Deze bandbreedte is de ondergrens om moderne administratieve 
    en technische dossiers binnen een acceptabel tijdstvenster te kunnen synchroniseren, 
    in tegenstelling tot de trage, puur op tekst gerichte legacy-systemen zoals Winlink 
    \autocite{WinlinkOverview}.
\end{itemize}

\subsection{Maximale latency (REQ-02)}
In tegenstelling tot UDP-gebaseerde mediastromen, maken bestandsoverdrachtsprotocollen 
veelal gebruik van TCP op de transportlaag voor betrouwbare aflevering \autocite{Tanenbaum2021}. 
TCP-venstergrootte-algoritmen (*sliding window*) zijn extreem gevoelig voor de 
*Round-Trip Time* (RTT).
\begin{itemize}
    \item \textbf{Eis:} De gemiddelde RTT (latency) over de radioverbinding mag niet hoger 
    zijn dan $150\text{ ms}$, met een maximale jitter (variantie in latency) van $\pm 20\text{ ms}$.
    \item \textbf{Verantwoording:} Boven deze drempelwaarde treedt er door de interactieve 
    aard van protocollen zoals SMB of SFTP een exponentiële degradatie van de effectieve 
    doorvoersnelheid op, onafhankelijk van de beschikbare rauwe bandbreedte \autocite{Pollefliet2011}.
\end{itemize}

\subsection{Packet Loss Tolerantie (REQ-03)}
Atmosferische storingen en multi-path fading veroorzaken bitfouten op de fysieke laag van 
draadloze verbindingen \autocite{Tanenbaum2021}. 
\begin{itemize}
    \item \textbf{Eis:} Het netwerk moet een packet loss-percentage van maximaal $2\%$ tolereren 
    tijdens continue bestandsoverdracht, zonder dat de onderliggende TCP-sessie onherstelbaar 
    afbreekt.
    \item \textbf{Verantwoording:} Wanneer packet loss deze grens overschrijdt, interpreteert het 
    TCP-congestie-algoritme dit als netwerkcongestie, waarna de verzendsnelheid (*congestion window*) 
    drastisch wordt gereduceerd tot een fractie van de kanaalcapaciteit.
\end{itemize}

\section{Wettelijke en Reglementaire Randvoorwaarden}
\label{sec:req-wettelijk}

De fysieke implementatie van de radioverbinding is gebonden aan dwingende nationale wetgeving 
om interferentie met primaire communicatiediensten uit te sluiten. Dit leidt tot harde 
randvoorwaarden voor de hardwareconfiguratie.

\subsection{Conformiteit met het BIPT-frequentieplan (REQ-04)}
\begin{itemize}
    \item \textbf{Eis:} De proof-of-concept mag uitsluitend opereren binnen de frequenties, 
    transmissiemodi en maximale vermogenslimieten (ERP/PEP) zoals wettelijk vastgelegd in het 
    \textcite{BIPT2019}.
    \item \textbf{Specificatie:} Bij gebruik van het radioamateurspectrum (bijvoorbeeld de 70~cm-band) 
    moet voldaan worden aan de specifieke bandplannen gecureerd door de Unie van Belgische 
    Zendamateurs \autocite{UBAFrequenties}. Indien wordt uitgeweken naar licentievrije 
    alternatieven (2.4~GHz of 5.6~GHz), moet de hardware strikt begrensd zijn op de wettelijke 
    EIRP-limieten voor ISM-toepassingen om sancties te vermijden.
\end{itemize}

\section{Tussentijdse Conclusie}
\label{sec:req-tussentijdse-conclusie}

De requirementsanalyse en de daaruit voortvloeiende spectrum- en protocolvergelijking 
leiden tot een eenduidige blauwdruk voor de te realiseren proof-of-concept. Om een 
degelijke, performante en infrastructuuronafhankelijke netwerkverbinding met een 
fileserver op te zetten, is de combinatie van een **5.6~GHz Point-to-Point straalverbinding** op de fysieke laag en het **SFTP-protocol** op de applicatielaag de meest optimale 
architectuur.

De 5.6~GHz-band levert via breedbandige OFDM-modulatie de noodzakelijke kanaalcapaciteit 
om consistent aan de hoge doorvoersnelheid (REQ-01) te voldoen, zonder de extreme 
atmosferische zuurstofabsorptie te ondergaan die de 60~GHz-band kenmerkt. Op het vlak van 
regelgeving kent deze band in België een belangrijk duaal karakter conform het 
\textcite{BIPT2019}:
\begin{itemize}
    \item \textbf{Licentievrij (RLAN/ISM):} Binnen het segment $5470\text{--}5725\text{ MHz}$ 
    is \textbf{geen radioamateurvergunning vereist}, mits het zendvermogen strikt begrensd 
    blijft op maximaal \textbf{1~Watt ($30\text{ dBm}$) EIRP} en de hardware verplicht gebruikmaakt 
    van *Dynamic Frequency Selection* (DFS) om interferentie met weerradars te voorkomen.
    \item \textbf{Gelicentieerd (Amateursubband):} Indien men gebruikmaakt van het secundaire 
    radioamateursegment ($5650\text{--}5850\text{ MHz}$) vervalt de DFS-verplichting en stijgt 
    het wettelijke vermogensprivilege tot maximaal \textbf{200~Watt PEP} (voor Klasse A/HAREC), 
    wat echter een strikte persoonlijke vergunningsplicht met zich meebrengt \autocite{UBAFrequenties}.
\end{itemize}
Voor de basisopstelling van de PoC wordt uitgegaan van de licentievrije limieten ($1\text{~W EIRP}$), 
zodat de oplossing universeel inzetbaar blijft.

Op de applicatielaag is SFTP onmisbaar gebleken: het omzeilt de *chattiness* en de daaruit 
voortvloeiende throughput-instorting van protocollen zoals SMB bij latency- en jitter-fluctuaties 
(REQ-02). Bovendien blijft SFTP stabiel functioneren binnen het gestelde packet loss-budget van 
maximaal 2\% (REQ-03), terwijl het direct de vereiste databeveiliging introduceert.

Wat betreft het operationele bereik is deze 5.6~GHz/SFTP-architectuur opgedeeld in de 
volgende specifieke afstandssegmenten:
\begin{itemize}
    \item \textbf{0 tot 15 kilometer (Optimaal bereik):} Binnen deze range functioneert de 
    architectuur maximaal. Zolang er een vrije optische zichtlijn (*Line-of-Sight*) is en de 
    eerste Fresnel-zone vrij is van obstructies, wordt de *high-performance* baseline van 
    $10\text{ Mbit/s}$ tot over de $100\text{ Mbit/s}$ stabiel behaald met minimale latency ($<10\text{ ms}$).
    \item \textbf{15 tot 30 kilometer (Horizonlimiet):} De verbinding blijft bruikbaar en 
    voldoet aan de minimale ondergrens van $1\text{ Mbit/s}$ (REQ-01), mits de antennes op 
    voldoende hoogte (masten, daken of heuvels) worden geplaatst om de fysieke kromming van de 
    aarde te compenseren. SFTP vangt hier de incidentele stijging in packet loss door vrije-ruimtedemping op.
    \item \textbf{Groter dan 30 kilometer (Suboptimaal/Extreme afstand):} Dit bereik valt buiten 
    het betrouwbare toepassingsgebied voor permanente bestandstoegang wegens atmosferische 
    onvoorspelbaarheid, waardoor hier multi-hop routing of een overstap naar het smalbandige UHF-spectrum noodzakelijk wordt.
\end{itemize}

\subsection{Impact van de Licentievrije Limieten op het Projectontwerp}
\label{sec:req-impact-licentievrij}

Het operationeel inzetten van de 5.6~GHz-band zonder licentie (binnen het $5470\text{--}5725\text{ MHz}$ RLAN-segment) legt een aantal dwingende wettelijke en technische restricties op. Deze limieten hebben een directe impact op de configuratie en de te verwachten netwerkprestaties van de proof-of-concept:

\begin{enumerate}
    \item \textbf{EIRP-Vermogenslimiet (Max. 1 Watt / $30\text{ dBm}$):} \\
    De absolute juridische grens voor het effectief uitgestraalde vermogen (EIRP) beperkt het totale *link-budget*. In de praktijk betekent dit dat bij de inzet van high-gain richtantennes (essentieel voor afstanden van $15\text{--}30\text{ km}$), het zendvermogen van de radiofrequentie-uitgang softwarematig moet worden teruggeschroefd om de 1~Watt-grens niet te overschrijden. Dit begrenst de signaalmarge bij slechte weersomstandigheden.
    
    \item \textbf{DFS-Verplichting (*Dynamic Frequency Selection*):} \\
    De hardware is wettelijk verplicht om passief te scannen naar prioritair spectrumgebruik, zoals weerradars. Indien een radarsignaal wordt gedetecteerd, dwingt de wet de straalverbinding om onmiddellijk van frequentie te wisselen. Hoewel het SFTP/TCP-protocol door zijn stateful karakter bestand is tegen de hieruit voortvloeiende korte onderbrekingen, introduceert dit incidentele latency-pieken en jitter (REQ-02) tijdens het heronderhandelen van de link.
    
    \item \textbf{Gedeeld Medium (*Shared Spectrum* en Congestie):} \\
    Omdat de band vrij toegankelijk is voor commerciële en private toepassingen, is er in dichtbevolkte gebieden geen garantie op een storingsvrij transmissiekanaal. Eventuele *co-channel interferentie* verhoogt de bitfoutfrequentie (BER) en daarmee het packet loss-percentage. Om de packet loss binnen het budget van maximaal 2\% (REQ-03) te houden, dwingt deze limiet het project tot het selecteren van smallere kanaalbreedtes (bijv. $20\text{ MHz}$ in plaats van $80\text{ MHz}$), wat de theoretische maximale doorvoersnelheid reduceert maar de link-robuustheid maximaliseert.
\end{enumerate}

Ter synthese toont de onderstaande matrix de directe correlatie tussen de wettelijke bepalingen en de netwerktechnische requirements van dit onderzoek:

\begin{table}[h]
    \centering
    \small
    \begin{tabular}{lll}
        \toprule
        \textbf{BIPT-Limiet (Licentievrij)} & \textbf{Technisch gevolg voor de PoC} & \textbf{Impact op Requirements} \\
        \midrule
        Max. 1~W EIRP & Beperkt de maximale signaalmarge & Maakt afstanden $>15\text{ km}$ topografisch kritiek \\
        DFS-Verplichting & Abrupte kanaalwissels bij radardetectie & Veroorzaakt tijdelijke latency-pieken (REQ-02) \\
        Gedeeld Spectrum & Risico op co-channel interferentie & Verhoogt packet loss (REQ-03); dwingt tot smalle kanalen \\
        \bottomrule
    \end{tabular}
    \caption{Synthese van licentievrije restricties en de impact op de netwerkrequirements.}
    \label{tab:limiet-impact-matrix}
\end{table}

Dit afgebakende theoretische, functionele en wettelijke kader minimaliseert de architecturale 
risico's. Het vormt de directe deliverablesleutel voor het volgende hoofdstuk, waarin de fysieke 
assemblage, de hardwareconfiguratie en het daadwerkelijke netwerkontwerp van de testopstelling 
worden gerealiseerd.
